\subsection{子浸入}

5.10.1. 设 \(f:X\to Y\) 是流形的一个态射，\(\Gamma\) 是 \(f\) 的图。映射 \(j:x\mapsto(x,f(x))\) 是 X 到 \(X\)\(X\times Y\) 的浸入，其像是子流形 \(\Gamma\)；而 \(f=\mathrm{pr}_{2}\circ j\) 是先浸入 \(j\)、再经浸没 \(\mathrm{pr}_{2}\) 所得的复合。

令 \(a \in X\)。若存在 a 的一个开邻域 、一个流形 、从 U 到 Z 的浸没  以及从 Z 到 Y 的浸入 ，使得 ，则称 \(f\) 在 \(a\) 处是子浸入。\(U\)\(a\)\(Z\)\(s\)\(U\)\(Z\)\(i\)\(Z\)\(Y\)\(f|U = i \circ s\)\(X\)\(f\) 为子浸入的 X 中点集是 X 的开子集；若这个开集就是 X 本身，则称 \(X\)\(X\)\(f\) 是子浸入。

5.10.2. 浸入和浸没都是子浸入。若 \(f: X \to Y\) 是子浸入，\(g: Y \to Z\) 是浸入，且 \(h: W \to X\) 是浸没，则 \(g \circ f \circ h\) 是子浸入。

若 \(f\) 和 \(f'\) 是子浸入，则 \(f \times f'\) 是子浸入。

5.10.3. \(f: X \to Y\) 在 X 中一点 \(a\) 处为子浸入的充分必要条件是：存在 \(X\) 在 a 处的一个图 \((U, \varphi, E)\)、Y 在点 \(X\)\(a\) 处的一个图 \((V, \psi, F)\)，以及从 E 到 F 的连续线性映射 \(Y\)\(f(a)\)\(g\)，使得：\(E\)\(F\)

\begin{enumerate}
\def\labelenumi{(\roman{enumi})}
\item
  \(f(\mathbf{U}) \subset \mathbf{V}\) , \(g(\varphi(\mathbf{U})) \subset \psi(\mathbf{V})\) 且 \(f|\mathbf{U} = \psi^{-1} \circ g \circ \varphi\) ;
\item
  g 的核（分别地，像）是 E（分别地，F）的具有拓扑补空间的闭子空间。
\end{enumerate}

5.10.4. 设 X 和 Y 是两个有限维流形。态射 \(f\) 从 X 到 Y 在 X 中一点 \(a\) 处为子浸入的充分必要条件是：存在 X 在 \((\xi^{1},\ldots,\xi^{m})\) 处的坐标系 \(a\)、Y 在 \((\eta^{1},\ldots,\eta^{n})\) 处的坐标系 \(f(a)\)，以及一个整数 \(k\leqslant\inf(m,n)\)，使得

\[
\eta^ {i} \circ f = \xi^ {i} \quad \text { for } 1 \leqslant i \leqslant k
\]

\[
\eta^ {i} \circ f = 0 \quad \text {   for   } k <   i \leqslant n.
\]

5.10.5. 设 \(f: X \to Y\) 是子浸入。对于每个 \(b \in Y\)，\(f^{-1}(b)\) 是 X 的子流形；在点 \(X\) 处与子流形  相切的 \(\mathrm{T}_{a}(X)\) 的子空间是 \(f^{-1}(b)\)\(a \in f^{-1}(b)\)\(\mathrm{T}_{a}(f)\) 的核。

5.10.6.（“常秩定理”）设 \(f: X \to Y\) 是流形的一个态射，且 \(a \in X\)。若 \(f\) 在 \(a\) 处是子浸入，则当 \(\operatorname{rg}_{x} f = \operatorname{rg}_{a} f\) 充分接近 \(x\) 时，有 \(a\)。

反之，假设域 K 的特征为 0。设 (U, \(\varphi\) , E) 是 X 在 \(a\) 处的一个图，(V, \(\psi\) , F) 是 Y 在 \(f(a)\) 处的一个图，且 \(f(U) \subset V\)。置 \(g = \psi \circ f \circ \varphi^{-1}\)。若存在 E 的一个闭向量子空间 \(E_{1}\) 和 F 的一个闭向量子空间 \(F_{1}\)，使得对于每个 \(x \in U\)，子空间 \(E_{1}\)（分别为 \(F_{1}\)）是 g 在点 \(\varphi(x)\) 处的导数 Dg( \(g\)\(\varphi(x)\) ) 的核（分别为像）的拓扑补空间，则 \(f\) 在 \(a\) 处是子浸入。

若 K 的特征为 0（分别地，K = R 或 C），且 Y 是有限维的（分别地，\(rg_{a}f < +\infty\)），则 f 在 a 处是子浸入，当且仅当对于所有充分接近 a 的 x，都有 \(rg_{x}f = rg_{a}f\)。

若 K 的特征为 0（分别地，K = R 或 C），且 Y（分别地，X）是有限维的，则 f 为子浸入的点集 \(x \in X\) 是 X 的稠密开子集。

5.10.7.（“子浸入的典范分解”）每个子浸入都是一个浸没与一个浸入的复合。具体地，设 \(f: X \to Y\) 是子浸入，J 是 Y 的子流形芽的流形（5.8.12）。令 \(x\) 属于 X 且 \(y = f(x)\)；存在 x 的开邻域 U，使得 \(x\)\(f|U\) 是从 U 到 Y 的一个子流形 Z 上的浸没；J 的元素 \(\gamma_y(Z)\) 只依赖于 \(x\)，而不依赖于所选的邻域 U，若将其记为 \(\lambda(x)\)，则映射 \(\lambda\) 是从 X 到 J 的浸没；若以 \(\rho\) 表示 J 到 Y 的典范浸入，则有 \(f = \rho \circ \lambda\)。若 \(f\) 是浸入，则从 X 到 J 的态射 \(\lambda\) 是 étale 的。

若 g 是从 X 到流形 Z 的满浸没，h 是从 Z 到 Y 的态射，且 \(f = h \circ g\)，则存在唯一的从 Z 到 J 的浸没 \(\mu\)，使得 \(\lambda = \mu \circ g\)。

